\section{Transmisi\'on, ingesta e interfaces}
\label{sec:transmision-interfaces}

Una vez que la se\~nal fue acondicionada y digitalizada, el problema deja de ser
exclusivamente anal\'ogico. Los datos deben abandonar cada nodo, llegar sin
ambig\"uedades a un punto com\'un, conservar su contexto de adquisici\'on y quedar
disponibles para revisi\'on y procesamiento. En esta etapa aparecen, por tanto,
dos interfaces distintas:

\begin{equation}
\text{nodos de medida}
\longrightarrow
\text{maestro de campo}
\longrightarrow
\text{servidor de ingesta}
\longrightarrow
\text{revisi\'on y MASW}.
\end{equation}

La primera interfaz permite controlar la adquisici\'on: armar nodos, ajustar
par\'ametros, iniciar o detener una captura, observar el enlace y exportar los
datos. La segunda permite controlar el conjunto de datos: validar capturas,
filtrar, agrupar, alinear, promediar, construir el \emph{gather}, estimar la
dispersi\'on e invertir el perfil de velocidad. Esta separaci\'on evita que el
microcontrolador tenga que ejecutar las tareas pesadas y permite que el trabajo
de campo contin\'ue aun cuando no exista conexi\'on a Internet.

\subsection{Decisi\'on de plataforma de comunicaci\'on}

Se seleccion\'o la familia ESP32 por accesibilidad local, bajo costo, abundante
documentaci\'on, integraci\'on de Wi-Fi y disponibilidad de una pila de desarrollo
simple. No se plantea al ESP32 como un instrumento de medida: su funci\'on comienza
despu\'es de la digitalizaci\'on. El nodo esclavo recibe los datos digitales del
frente de adquisici\'on, los empaqueta y los entrega al maestro.

Para el enlace entre nodos se adopt\'o ESP-NOW. Esta decisi\'on permite una red
directa de baja sobrecarga que no depende de un \emph{router} en el sitio. A
cambio, impone restricciones concretas: carga \`util limitada por paquete,
entrega no equivalente a un transporte confiable, necesidad de mantener el canal
de radio y dificultad para sostener simult\'aneamente la red de adquisici\'on y
una conexi\'on convencional a Internet. El dise\~no afronta estas limitaciones
mediante fragmentaci\'on y reensamblado, numeraci\'on de mensajes, confirmaciones,
reintentos, estados expl\'icitos de cada nodo y almacenamiento local antes de la
subida.

El maestro concentra los enlaces ESP-NOW y expone una red Wi-Fi propia. Desde un
tel\'efono o una computadora se abre la interfaz web embebida, sin instalar una
aplicaci\'on. Cuando hay conectividad disponible, el equipo cambia a la fase de
enlace, se asocia a una red de \SI{2.4}{\giga\hertz} y ofrece la cola de capturas
para su transferencia. El uso de un nombre estable mediante mDNS evita que el
operador tenga que memorizar la direcci\'on que recibe el maestro en cada red.

\subsection{Por qu\'e se desarrollaron varias interfaces}

Las interfaces no fueron implementaciones redundantes de una misma idea. Cada
una resolvi\'o una pregunta de una etapa diferente y dej\'o funciones que luego
se trasladaron a la siguiente versi\'on.

\begin{table}[ht]
\centering
\small
\begin{tabularx}{\textwidth}{@{}p{0.20\textwidth}p{0.23\textwidth}Y@{}}
\toprule
\textbf{Interfaz} & \textbf{Uso principal} & \textbf{Aporte al dise\~no} \\
\midrule
MATLAB: alcance USB & Verificar el canal individual & Visualizaci\'on cruda, referencia DC, FIR, persistencia de configuraci\'on y guardado de muestras. \\
MATLAB: InterfaceESP & Probar maestro y esclavos & Sincronizaci\'on, lotes, canales, control de PGA, FIR, VDAC y depuraci\'on por nodo. \\
PyQt: adquisici\'on & Operar varios nodos & Escritorio multipesta\~na, estado de enlace, gr\'aficos simult\'aneos, tema oscuro/claro y registro central. \\
SPA embebida & Operar en campo sin instalar software & Control desde cualquier navegador, funcionamiento fuera de l\'inea y configuraci\'on del enlace. \\
PyQt: revisi\'on & Explorar el flujo cient\'ifico & Prototipo de validaci\'on, filtros, grupos, enfase, promedios, \emph{waterfall} y MASW. \\
FastAPI web & Ingestar y procesar en una direcci\'on conocida & Acceso remoto, estado persistente y continuidad entre campa\~nas y equipos. \\
\bottomrule
\end{tabularx}
\caption{Funciones y alcance de las interfaces desarrolladas.}
\label{tab:evolucion-interfaces}
\end{table}

\clearpage
\subsection{Primera interfaz MATLAB: observaci\'on del canal}

La primera aplicaci\'on cumple la funci\'on de osciloscopio de desarrollo. Su
importancia no reside en la apariencia, sino en que concentra las variables que
deben quedar controladas antes de comparar capturas: puerto USB, referencia DC,
ganancia, filtro, cantidad de muestras, escala y persistencia de los bloques
anal\'ogicos.

\begin{figure}[p]
\centering
\includegraphics[width=0.88\textwidth,height=0.82\textheight,keepaspectratio]{matlab_scope_simple.png}
\caption{Interfaz MATLAB de alcance individual. Se distinguen la conexi\'on USB,
la referencia DC, el FIR calculado en MATLAB, el control de PGA, el guardado de
muestras y los controles de escala y desplazamiento. La pantalla permite
contrastar la forma de onda cruda con el resultado filtrado antes de introducir
la comunicaci\'on distribuida.}
\label{fig:matlab-scope-simple}
\end{figure}

\clearpage
\subsection{MATLAB InterfaceESP: del canal aislado al sistema distribuido}

La segunda aplicaci\'on incorpora al maestro y a los esclavos. La vista
\textbf{Stream} concentra el estado global, mientras que cada pesta\~na de
esclavo expone sus par\'ametros. Esta organizaci\'on hace visible una decisi\'on
central: los nodos comparten una l\'ogica com\'un, pero deben conservar identidad,
estado y diagn\'ostico propios.

\begin{figure}[p]
\centering
\includegraphics[width=\textwidth,height=0.80\textheight,keepaspectratio]{matlab_interface_esp_05_stream.png}
\caption{Vista \textbf{Stream} de InterfaceESP. Incluye puerto serie, cantidad
de esclavos, sincronizaci\'on, recepci\'on por lotes, respuesta de ESP-NOW,
selecci\'on de canal y guardado. La interfaz ya trata la adquisici\'on como un
sistema de nodos y no como un osciloscopio aislado.}
\label{fig:matlab-interface-stream}
\end{figure}

\begin{figure}[p]
\centering
\includegraphics[width=\textwidth,height=0.80\textheight,keepaspectratio]{matlab_interface_esp_04_esclavo_1_s1.png}
\caption{Pesta\~na del esclavo 1. Se muestran los controles de VDAC, PGA, FIR,
se\~nal de prueba, cancelaci\'on de \SI{50}{\hertz} y puerto de depuraci\'on. La
separaci\'on por nodo permite calibrar o diagnosticar uno sin ocultar el estado
global.}
\label{fig:matlab-interface-s1}
\end{figure}

\begin{figure}[p]
\centering
\includegraphics[width=\textwidth,height=0.80\textheight,keepaspectratio]{matlab_interface_esp_03_esclavo_2_s2.png}
\caption{Pesta\~na del esclavo 2, generada con una configuraci\'on de tres nodos.
La repetici\'on controlada del panel comprueba que la aplicaci\'on no depend\'ia de
un \`unico esclavo codificado de forma fija.}
\label{fig:matlab-interface-s2}
\end{figure}

\begin{figure}[p]
\centering
\includegraphics[width=\textwidth,height=0.80\textheight,keepaspectratio]{matlab_interface_esp_02_esclavo_3_s3.png}
\caption{Pesta\~na del esclavo 3. La identidad del nodo forma parte del estado de
la interfaz y debe acompa\~nar a la captura hasta el servidor.}
\label{fig:matlab-interface-s3}
\end{figure}

\begin{figure}[p]
\centering
\includegraphics[width=\textwidth,height=0.80\textheight,keepaspectratio]{matlab_interface_esp_01_log.png}
\caption{Registro de InterfaceESP. El \emph{log} desacopla el diagn\'ostico de la
visualizaci\'on: conexiones, confirmaciones, errores y transiciones pueden
auditarse sin depender de que una curva se vea correctamente.}
\label{fig:matlab-interface-log}
\end{figure}

\clearpage
\subsection{Aplicaci\'on PyQt de adquisici\'on multipunto}

La versi\'on PyQt reorganiza el prototipo en una herramienta de escritorio. La
vista global presenta las trazas de varios nodos en paralelo y mantiene acceso
directo a cada esclavo y al registro. Tambi\'en se ensayaron temas oscuro y
claro, relevante para una herramienta que puede utilizarse tanto en laboratorio
como a plena luz durante una campa\~na.

\begin{figure}[p]
\centering
\includegraphics[width=\textwidth,height=0.80\textheight,keepaspectratio]{qt_scope_01_stream.png}
\caption{Vista global oscura de \textbf{Geophone Multi-Node Scope}. El panel
superior concentra conexi\'on y acciones globales; el resto del espacio se
reserva a la comparaci\'on temporal entre canales.}
\label{fig:qt-scope-stream}
\end{figure}

\begin{figure}[p]
\centering
\includegraphics[width=\textwidth,height=0.80\textheight,keepaspectratio]{qt_scope_02_esclavo1.png}
\caption{Controles del esclavo 1 en la aplicaci\'on PyQt. La configuraci\'on y el
diagn\'ostico se agrupan por nodo sin abandonar la misma sesi\'on de adquisici\'on.}
\label{fig:qt-scope-s1}
\end{figure}

\begin{figure}[p]
\centering
\includegraphics[width=\textwidth,height=0.80\textheight,keepaspectratio]{qt_scope_03_esclavo2.png}
\caption{Controles del esclavo 2. La estructura equivalente facilita detectar
diferencias de configuraci\'on que podr\'ian confundirse con diferencias del
subsuelo.}
\label{fig:qt-scope-s2}
\end{figure}

\begin{figure}[p]
\centering
\includegraphics[width=\textwidth,height=0.80\textheight,keepaspectratio]{qt_scope_04_esclavo3.png}
\caption{Controles del esclavo 3. Esta vista confirma la extensi\'on del patr\'on
de operaci\'on a tres unidades de adquisici\'on.}
\label{fig:qt-scope-s3}
\end{figure}

\begin{figure}[p]
\centering
\includegraphics[width=\textwidth,height=0.80\textheight,keepaspectratio]{qt_scope_05_log.png}
\caption{Registro central de la aplicaci\'on PyQt. La secuencia de eventos de
comunicaci\'on permanece disponible aunque el operador cambie de pesta\~na.}
\label{fig:qt-scope-log}
\end{figure}

\begin{figure}[p]
\centering
\includegraphics[width=\textwidth,height=0.80\textheight,keepaspectratio]{qt_scope_06_stream_tema_claro.png}
\caption{La misma vista global en tema claro. Esta variante permite evaluar la
legibilidad de trazas, controles y estados bajo otra condici\'on de iluminaci\'on.}
\label{fig:qt-scope-light}
\end{figure}

\clearpage
\subsection{Interfaz embebida en el maestro ESP32}

El paso siguiente consiste en eliminar la dependencia de MATLAB, Python o un
instalador en la computadora de campo. La aplicaci\'on se almacena junto al
firmware y se entrega como una SPA desde el propio maestro. Su historial Git
permite reconstruir siete hitos de desarrollo, todos ejecutables sin hardware
mediante un servidor que emula las rutas del ESP32.

\begin{table}[ht]
\centering
\small
\begin{tabularx}{\textwidth}{@{}p{0.13\textwidth}p{0.15\textwidth}Y@{}}
\toprule
\textbf{Fecha} & \textbf{Commit} & \textbf{Cambio visible o funcional} \\
\midrule
08/06/2026 & \texttt{323fcb7} & Primera interfaz web, relay WebSocket y pruebas. \\
08/06/2026 & \texttt{76d6092} & Primer estado funcional, todav\'ia incompleto. \\
15/06/2026 & \texttt{70ca42e} & Incorporaci\'on de la autocalibraci\'on. \\
01/07/2026 & \texttt{8b63d2b} & P\'agina operativa para el flujo de adquisici\'on. \\
12/07/2026 & \texttt{89384d8} & Cierre de la UI y herramientas de captura de campo. \\
24/07/2026 & \texttt{9eb8b5a} & Configuraci\'on de enlace, escaneo de redes, canal y mDNS. \\
03/08/2026 & \texttt{d6463a} & Compilaci\'on local estable y selecci\'on PGA/PGAout. \\
\bottomrule
\end{tabularx}
\caption{Versiones de la SPA extra\'idas del historial del firmware ESP32.}
\label{tab:spa-historica}
\end{table}

La inspecci\'on de los documentos HTML de cada commit permite reconstruir
tambi\'en la evoluci\'on de los men\'us. El cambio no fue solamente visual: a
medida que la adquisici\'on dej\'o de ser una prueba de WebSocket aparecieron
operaciones de campo que antes no exist\'ian.

\begin{table}[ht]
\centering
\scriptsize
\begin{tabularx}{\textwidth}{@{}p{0.11\textwidth}p{0.25\textwidth}Y@{}}
\toprule
\textbf{Hito} & \textbf{Men\'us visibles} & \textbf{Controles incorporados o consolidados} \\
\midrule
08/06 inicial & Vista \`unica & Reconexi\'on, tema, ARM, START/STOP, estado, zoom, reset y descarga ZIP. \\
08/06 funcional & Captura, Nodos, Esclavos, Export, Log & Separaci\'on entre operaci\'on global, estado por nodo, exportaci\'on y diagn\'ostico. \\
15/06 & Captura, Nodos, Esclavos, Export, Log & Cursores, espectro y autocalibraci\'on; la interfaz comienza a asistir la evaluaci\'on de la se\~nal. \\
01/07 & Captura, Esclavos, Preservados, Log & Preservaci\'on expl\'icita, calibraci\'on de todos los nodos, guardado CSV/ZIP y autenticaci\'on. \\
12/07 & Captura, Esclavos, Preservados, Log & Decimaci\'on por $N$, duraciones separadas, l\'imites de RAM/SD y operaci\'on de descuento. \\
24/07 & Captura, Esclavos, Preservados, Enlace, Log & Escaneo de redes, credenciales, URL de ingesta, cola de archivos, subida al servidor y mDNS. \\
03/08 & Captura, Esclavos, Preservados, Enlace, Log & Se estabiliza el flujo anterior y se incorpora la selecci\'on entre PGA y PGAout en la cadena distribuida. \\
\bottomrule
\end{tabularx}
\caption{Evoluci\'on verificable de men\'us y controles de la interfaz embebida.}
\label{tab:spa-menus}
\end{table}

La versi\'on reciente distribuye sus funciones en cinco men\'us:

\begin{itemize}
  \item \textbf{Captura}: armado, duraci\'on de ge\'ofono y martillo,
  decimaci\'on, l\'imites de RAM y SD, inicio, detenci\'on, calibraci\'on,
  exportaci\'on ZIP/CSV y subida al servidor.
  \item \textbf{Esclavos}: estado y configuraci\'on individual de cada nodo.
  \item \textbf{Preservados}: inspecci\'on de adquisiciones conservadas antes de
  su exportaci\'on o eliminaci\'on.
  \item \textbf{Enlace}: escaneo de redes, SSID, contrase\~na, URL del servidor,
  canal y estado de la cola.
  \item \textbf{Log}: secuencia de conexi\'on, comandos, confirmaciones y errores.
\end{itemize}

El gr\'afico global puede alternar se\~nal cruda, filtrada, envolvente,
componente DC y espectro; adem\'as ofrece zoom, cursores y tema. Esta densidad de
funciones es posible porque el navegador ejecuta la presentaci\'on y parte del
procesamiento visual, mientras el ESP32 conserva las tareas deterministas de
adquisici\'on, enlace y almacenamiento.

Las Tablas~\ref{tab:spa-historica} y~\ref{tab:spa-menus} ya documentan, hito a
hito, qu\'e problema resolvi\'o cada revisi\'on: de una p\'agina \'unica de prueba de
WebSocket a un flujo con captura, esclavos, preservados, enlace y registro. Las
dos figuras siguientes bastan para hacer visible ese recorrido dentro del texto
principal ---el punto de partida y el punto de llegada---; el resto de las
capturas de los siete hitos, incluidas las pesta\~nas de configuraci\'on repetidas
por nodo, se conserva de forma exhaustiva en el
Ap\'endice~\ref{sec:anexo-spa-visual} como banco visual de referencia, sin
interrumpir aqu\'i la lectura del flujo de transmisi\'on.

\begin{figure}[p]
\centering
\includegraphics[width=0.88\textwidth,height=0.82\textheight,keepaspectratio]{esp_historico/01_ui_inicial/01_vista_inicial_oscura.png}
\caption{Hito 1 (08/06/2026), p\'agina \'unica: adquisici\'on, gr\'afica, configuraci\'on
de esclavos y registro conviven todav\'ia sin separaci\'on por pesta\~nas. Esta es la
primera versi\'on ejecutable de la SPA, apenas posterior a que los datos empezaran a
salir del PSoC por I$^2$C.}
\label{fig:spa-hito1-resumen}
\end{figure}

\begin{figure}[p]
\centering
\includegraphics[width=0.88\textwidth,height=0.82\textheight,keepaspectratio]{esp_historico/07_pga_pgaout/01_captura_completa.png}
\caption{Hito 7 (03/08/2026), estado vigente: cinco men\'us (Captura, Esclavos,
Preservados, Enlace, Log), decimaci\'on, l\'imites de RAM/SD, selecci\'on PGA/PGAout y
subida al servidor. La comparaci\'on directa con la Fig.~\ref{fig:spa-hito1-resumen}
resume en dos im\'agenes lo que las Tablas~\ref{tab:spa-historica}
y~\ref{tab:spa-menus} describen en texto: la interfaz pas\'o de exponer transporte a
exponer una operaci\'on de campo completa.}
\label{fig:spa-hito7-resumen}
\end{figure}

\clearpage
\subsection{Revisi\'on de campo: prototipo PyQt del flujo cient\'ifico}

La herramienta de revisi\'on organiza el posprocesado como una secuencia
expl\'icita. Las capturas siguientes utilizan una copia de un conjunto real de
nueve adquisiciones; de ese modo la interfaz puede recorrerse sin modificar los
datos originales.

\begin{figure}[p]
\centering
\includegraphics[width=\textwidth,height=0.80\textheight,keepaspectratio]{qt_revision_01_capturas.png}
\caption{Pesta\~na \textbf{Capturas} de la revisi\'on PyQt. La tabla vincula cada
registro con las curvas de martillo y ge\'ofono, permitiendo inspeccionar la
forma de onda antes de cualquier operaci\'on colectiva.}
\label{fig:qt-review-capturas}
\end{figure}

\begin{figure}[p]
\centering
\includegraphics[width=\textwidth,height=0.80\textheight,keepaspectratio]{qt_revision_02_filtros.png}
\caption{Pesta\~na \textbf{Filtros}. Re\'une los par\'ametros que determinan qu\'e
contenido espectral entra a las etapas posteriores. Mantenerlos visibles evita
que el procesamiento se convierta en una operaci\'on oculta.}
\label{fig:qt-review-filtros}
\end{figure}

\begin{figure}[p]
\centering
\includegraphics[width=\textwidth,height=0.80\textheight,keepaspectratio]{qt_revision_03_agrupamiento.png}
\caption{Pesta\~na \textbf{Agrupamiento}. Las capturas se organizan por campa\~na,
distancia y criterio de repetici\'on antes de combinarse.}
\label{fig:qt-review-grupos}
\end{figure}

\begin{figure}[p]
\centering
\includegraphics[width=\textwidth,height=0.80\textheight,keepaspectratio]{qt_revision_04_enfase.png}
\caption{Pesta\~na \textbf{Enfase}. El alineamiento temporal utiliza la referencia
del golpe para que el promedio refuerce contenido coherente en lugar de borrar
los arribos por desfase.}
\label{fig:qt-review-enfase}
\end{figure}

\begin{figure}[p]
\centering
\includegraphics[width=\textwidth,height=0.80\textheight,keepaspectratio]{qt_revision_05_promedios.png}
\caption{Pesta\~na \textbf{Promedios}. Presenta los resultados combinados y los
tiempos de arribo derivados de grupos ya validados y alineados.}
\label{fig:qt-review-promedios}
\end{figure}

\begin{figure}[p]
\centering
\includegraphics[width=\textwidth,height=0.80\textheight,keepaspectratio]{qt_revision_06_waterfall.png}
\caption{Pesta\~na \textbf{Waterfall}. El conjunto de trazas se ordena por
distancia para observar el cambio temporal y construir el registro multicanal
equivalente que requiere MASW.}
\label{fig:qt-review-waterfall}
\end{figure}

\begin{figure}[p]
\centering
\includegraphics[width=\textwidth,height=0.80\textheight,keepaspectratio]{qt_revision_07_masw_dispersion.png}
\caption{Subpesta\~na \textbf{Dispersi\'on}. Los controles definen los rangos de
frecuencia y velocidad de fase y preparan la selecci\'on de la curva aparente
$V_R(f)$.}
\label{fig:qt-review-masw-disp}
\end{figure}

\begin{figure}[p]
\centering
\includegraphics[width=\textwidth,height=0.80\textheight,keepaspectratio]{qt_revision_08_masw_inversion.png}
\caption{Subpesta\~na \textbf{Inversi\'on}. La curva de dispersi\'on elegida se
convierte en entrada de un modelo estratificado cuyas velocidades y espesores
se ajustan para reducir el error entre dispersi\'on observada y calculada.}
\label{fig:qt-review-masw-inv}
\end{figure}

\begin{figure}[p]
\centering
\includegraphics[width=\textwidth,height=0.80\textheight,keepaspectratio]{qt_revision_09_masw_perfil_vs.png}
\caption{Subpesta\~na \textbf{Perfil Vs}. Es la salida final prevista del flujo:
un modelo de velocidad de onda de corte en funci\'on de la profundidad. La vista
permanece vac\'ia hasta ejecutar y aceptar una inversi\'on.}
\label{fig:qt-review-masw-vs}
\end{figure}

\clearpage
\subsection{Servidor web: ingesta y procesamiento desde una direcci\'on conocida}

La experiencia con las interfaces anteriores conduce a un modelo web para el
trabajo posterior a la adquisici\'on. El operador no necesita copiar una
configuraci\'on de MATLAB o Python entre computadoras y el estado del conjunto
puede persistir junto a los datos. El servidor se ejecuta en modo de solo
lectura para estas capturas y utiliza la misma copia de nueve adquisiciones.

\begin{figure}[p]
\centering
\includegraphics[width=\textwidth,height=0.80\textheight,keepaspectratio]{web_servidor_01_capturas_datos.png}
\caption{Servidor web, pesta\~na \textbf{Capturas}. La lista de adquisiciones se
relaciona con la se\~nal real de martillo y ge\'ofono. Esta es la puerta de
entrada para validar metadatos, polaridad, distancia y zona \`util.}
\label{fig:web-capturas}
\end{figure}

\begin{figure}[p]
\centering
\includegraphics[width=\textwidth,height=0.80\textheight,keepaspectratio]{web_servidor_02_filtros_kalman.png}
\caption{Servidor web, pesta\~na \textbf{Filtros}. Se muestran las alternativas
de preprocesado y el panel de configuraci\'on del filtro de Kalman.}
\label{fig:web-filtros}
\end{figure}

\begin{figure}[p]
\centering
\includegraphics[width=\textwidth,height=0.88\textheight,keepaspectratio]{web_servidor_03_filtros_kalman_expandido.png}
\caption{Configuraci\'on expandida del filtro de Kalman. La captura larga conserva
todos los campos y explicaciones del men\'u, aunque en una versi\'on resumida de
la presentaci\'on podr\'ia reemplazarse por un recorte.}
\label{fig:web-kalman-expandido}
\end{figure}

\begin{figure}[p]
\centering
\includegraphics[width=\textwidth,height=0.80\textheight,keepaspectratio]{web_servidor_04_agrupamiento.png}
\caption{Servidor web, pesta\~na \textbf{Agrupamiento}. Define qu\'e repeticiones
representan una misma posici\'on y pueden participar en las operaciones
colectivas.}
\label{fig:web-agrupamiento}
\end{figure}

\begin{figure}[p]
\centering
\includegraphics[width=\textwidth,height=0.80\textheight,keepaspectratio]{web_servidor_05_enfase.png}
\caption{Servidor web, pesta\~na \textbf{Enfase}. La interfaz hace expl\'icita la
referencia utilizada para alinear los golpes antes del promedio.}
\label{fig:web-enfase}
\end{figure}

\begin{figure}[p]
\centering
\includegraphics[width=\textwidth,height=0.80\textheight,keepaspectratio]{web_servidor_06_promedios_arrivals.png}
\caption{Servidor web, pesta\~na \textbf{Promedios/arrivals}. Los resultados
promediados y los arribos se inspeccionan antes de construir la representaci\'on
por distancia.}
\label{fig:web-promedios}
\end{figure}

\begin{figure}[p]
\centering
\includegraphics[width=\textwidth,height=0.80\textheight,keepaspectratio]{web_servidor_07_waterfall.png}
\caption{Servidor web, pesta\~na \textbf{Waterfall}. La visualizaci\'on integra
trazas de distintas posiciones en el \emph{gather} que alimenta el an\'alisis de
ondas superficiales.}
\label{fig:web-waterfall}
\end{figure}

\begin{figure}[p]
\centering
\includegraphics[width=\textwidth,height=0.80\textheight,keepaspectratio]{web_servidor_08_masw_dispersion.png}
\caption{Servidor web, MASW: \textbf{Dispersi\'on}. La interfaz contiene los
rangos de c\'alculo, normalizaci\'on, pesos, selecci\'on modal y herramientas de
\emph{picking}.}
\label{fig:web-masw-disp}
\end{figure}

\begin{figure}[p]
\centering
\includegraphics[width=\textwidth,height=0.80\textheight,keepaspectratio]{web_servidor_09_masw_inversion.png}
\caption{Servidor web, MASW: \textbf{Inversi\'on}. Esta pantalla concentra la
parametrizaci\'on del modelo y la ejecuci\'on del ajuste entre curva observada y
curva te\'orica.}
\label{fig:web-masw-inv}
\end{figure}

\begin{figure}[p]
\centering
\includegraphics[width=\textwidth,height=0.80\textheight,keepaspectratio]{web_servidor_10_masw_perfil_vs.png}
\caption{Servidor web, MASW: \textbf{Perfil Vs}. La salida se reserva para el
modelo final y sus indicadores, una vez que se complete y acepte la inversi\'on.}
\label{fig:web-masw-vs}
\end{figure}

\begin{figure}[p]
\centering
\includegraphics[width=\textwidth,height=0.80\textheight,keepaspectratio]{web_servidor_12_borrado.png}
\caption{Pesta\~na \textbf{Borrado}. Las operaciones destructivas quedan
separadas del flujo de an\'alisis. Esta decisi\'on de interfaz reduce el riesgo
de confundir una exclusi\'on metodol\'ogica con la eliminaci\'on f\'isica de una
captura.}
\label{fig:web-borrado}
\end{figure}

\clearpage
\subsection{Resultados reproducibles fuera de la interfaz: estado actual y l\'imites}

Las pantallas anteriores documentan \emph{d\'onde} se ejecuta cada etapa. Para no
confundir una interfaz implementada con un resultado cient\'ifico validado, esta
subsecci\'on conserva tambi\'en las salidas num\'ericas generadas por los scripts de
an\'alisis. Provienen de una reconstrucci\'on de 21 trazas reales entre 10 y
\SI{50}{\meter}, muestreadas a \SI{1020}{\hertz}. El repositorio almacena la salida
el\'ectrica de la cadena de aceleraci\'on en voltios; por eso las comparaciones
espectrales pueden defenderse, pero no deben presentarse como una calibraci\'on absoluta
de aceleraci\'on del suelo.

\begin{figure}[p]
\centering
\includegraphics[width=\textwidth,height=0.82\textheight,keepaspectratio]{../../../src/interfaces/python/geophone_scope/kalman_deconv/reports/real_canchita_group1/01_traza_real_y_espectro.png}
\caption{Traza real a \SI{30}{\meter} y distribuci\'on espectral normalizada.
Se comparan la salida medida, el filtrado Butterworth SOS aplicado en ambos
sentidos y la reconstrucci\'on RTS. La normalizaci\'on por \'area entre 1 y
\SI{200}{\hertz} permite comparar contenido espectral aun cuando las curvas no
comparten una calibraci\'on absoluta de amplitud.}
\label{fig:resultado-real-traza-espectro}
\end{figure}

\begin{figure}[p]
\centering
\includegraphics[width=\textwidth,height=0.82\textheight,keepaspectratio]{../../../src/interfaces/python/geophone_scope/kalman_deconv/reports/real_canchita_group1/02_dispersion_rts_real.png}
\caption{Imagen de dispersi\'on calculada con las trazas reconstruidas por RTS.
La referencia externa en cian se utiliza solamente para evaluar; no inicializa
ni ajusta el seguimiento. La curva verde corresponde al \emph{tracker}
Kalman+RTS en frecuencia y los puntos blancos al selector autom\'atico anterior.}
\label{fig:resultado-dispersion-rts}
\end{figure}

\begin{figure}[p]
\centering
\includegraphics[width=\textwidth,height=0.82\textheight,keepaspectratio]{../../../src/interfaces/python/geophone_scope/kalman_deconv/reports/real_canchita_group1/02_dispersion_sos_filtfilt.png}
\caption{Imagen de dispersi\'on obtenida con el brazo de control SOS
\emph{filtfilt}. Al aplicar el mismo \emph{tracker} de frecuencia aparece una
curva semejante a la del brazo RTS. Esta comparaci\'on separa el efecto del
preprocesamiento temporal del efecto del algoritmo de selecci\'on de cresta.}
\label{fig:resultado-dispersion-sos}
\end{figure}

\begin{figure}[p]
\centering
\includegraphics[width=\textwidth,height=0.82\textheight,keepaspectratio]{../../../src/interfaces/python/geophone_scope/kalman_deconv/reports/real_canchita_group1/04_comparacion_picking.png}
\caption{Comparaci\'on contra la referencia externa entre el selector autom\'atico
anterior y Kalman+RTS. En la banda 8--\SI{29.83}{\hertz}, el error RMSE pasa de
138,2 a \SI{16.2}{\meter\per\second} para SOS y de 148,2 a
\SI{16.4}{\meter\per\second} para RTS. Es evidencia de mejora del
\emph{picking} en este conjunto, no una validaci\'on general ni metrol\'ogica del
sistema.}
\label{fig:resultado-comparacion-picking}
\end{figure}

\begin{figure}[p]
\centering
\includegraphics[width=\textwidth,height=0.82\textheight,keepaspectratio]{../../../src/interfaces/python/geophone_scope/kalman_deconv/reports/acceleration_kalman_inversion_2026-08-17/01_masw_acceleration_kalman_curve.png}
\caption{Imagen MASW y curva seleccionada para la inversi\'on preliminar. Se
utilizan 25 puntos entre 8 y \SI{21.71}{\hertz}; el l\'imite superior se impone
antes de que el espaciamiento de \SI{2}{\meter} obligue al seguimiento a saltar
a otra cresta. La referencia externa se superpone solamente para evaluaci\'on.}
\label{fig:resultado-masw-curva-inversion}
\end{figure}

\begin{figure}[p]
\centering
\includegraphics[width=\textwidth,height=0.82\textheight,keepaspectratio]{../../../src/interfaces/python/geophone_scope/kalman_deconv/reports/acceleration_kalman_inversion_2026-08-17/02_inversion_vs_profile.png}
\caption{Ajuste de la curva de Rayleigh y dos parametrizaciones de $V_S(z)$.
El modelo mon\'otono obtiene un RMSE de \SI{0.463}{\meter\per\second} respecto
del \emph{pick}; permitir dos inversiones reduce el residuo, pero introduce
capas oscilatorias sensibles a la parametrizaci\'on. La banda disponible resuelve
aproximadamente hasta \SI{5.4}{\meter}: las interfaces y el semiespacio no deben
interpretarse a\'un como un perfil geol\'ogico definitivo.}
\label{fig:resultado-inversion-vs-preliminar}
\end{figure}

\noindent Estas salidas completan la cadena demostrada

\begin{center}
\small
\fbox{trazas reales}
$\rightarrow$
\fbox{preprocesado}
$\rightarrow$
\fbox{dispersi\'on}
$\rightarrow$
\fbox{selecci\'on $V_R(f)$}
$\rightarrow$
\fbox{inversi\'on}
$\rightarrow$
\fbox{$V_S(z)$ preliminar}.
\end{center}

Lo que queda por probar no es la existencia de estas etapas, sino su validez fuera del
conjunto mostrado: repetibilidad entre grupos y campa\~nas, sensibilidad a la geometr\'ia
y a las m\'ascaras, incertidumbre del \emph{picking}, no unicidad de la inversi\'on y
comparaci\'on con informaci\'on independiente del sitio.

\clearpage
\subsection{Acceso remoto, VPN y tunelamiento}

El equipo de campo puede operar sin Internet, pero la ingesta remota necesita
una direcci\'on que el maestro o el navegador puedan alcanzar. Una IP privada
asignada por una red Wi-Fi no cumple esta condici\'on: cambia entre sitios y no
es enrutable desde fuera. La soluci\'on adoptada combina una VPN de malla y
tunelamiento para presentar el servidor bajo una direcci\'on conocida.

Esta capa no sustituye el almacenamiento local. La secuencia robusta es:

\begin{enumerate}
  \item adquirir y verificar la captura en la red local del maestro;
  \item preservar el archivo y sus metadatos en el equipo de campo;
  \item entrar en la fase de enlace cuando exista una red adecuada;
  \item transferir a la direcci\'on conocida y comprobar la respuesta de ingesta;
  \item conservar o liberar la copia local seg\'un la pol\'itica de respaldo.
\end{enumerate}

La VPN resuelve identidad y alcance entre equipos; el t\'unel publica el servicio
sin exigir configuraci\'on manual de puertos en cada red visitada. La URL queda
guardada en la pesta\~na \textbf{Enlace}, por lo que el operador no necesita
editar el firmware al cambiar el servidor.

\subsection{Conclusi\'on de la etapa}

La evoluci\'on MATLAB $\rightarrow$ PyQt $\rightarrow$ SPA embebida y servidor
web no representa un cambio meramente est\'etico. Separa responsabilidades:
adquisici\'on local tolerante a desconexiones, transporte verificable, ingesta
persistente y procesamiento reproducible. Las capturas muestran que cada bloque
del flujo posee una interfaz observable y que las operaciones metodol\'ogicas no
quedan escondidas en un script \`unico.

Quedan por completar pruebas integrales con todos los nodos conectados, medidas
de p\'erdida y latencia bajo interferencia, verificaci\'on de recuperaci\'on tras
cortes de energ\'ia o red, y cierre del ciclo de inversi\'on hasta producir un
perfil Vs aceptado. Estas pruebas pertenecen a la validaci\'on del sistema; no
alteran la separaci\'on funcional definida en esta secci\'on.

\ifdefined\cuartaresumen\else\def\cuartaresumen{0}\fi

\clearpage
\ifnum\cuartaresumen=1

\section*{Ap\'endice: un hito por revisi\'on de la SPA embebida}
\label{sec:anexo-spa-visual}
\addcontentsline{toc}{section}{Ap\'endice: un hito por revisi\'on de la SPA embebida}

Este ap\'endice complementa la Fig.~\ref{fig:spa-hito1-resumen} y la
Fig.~\ref{fig:spa-hito7-resumen} con una captura por hito ---no una por pesta\~na
ni por nodo---, de modo que se vea \'unicamente qu\'e cambi\'o entre revisiones
consecutivas. Las pesta\~nas de configuraci\'on repetidas por nodo, los temas
claro/oscuro y las variantes de cursores o espectro se omiten aqu\'i porque no
aportan una diferencia funcional nueva frente al hito anterior; esa versi\'on
exhaustiva se conserva en \texttt{cuarta-seccion.tex} para quien necesite el
banco visual completo.

El servidor de demostraci\'on emula WebSocket, estados, tres nodos, escaneo de
redes e ingesta local. Las formas de onda son \textbf{se\~nales sint\'eticas
utilizadas para activar y documentar la interfaz}; no se presentan como
mediciones de campo ni como validaci\'on del sistema. Cada figura muestra una
sola pantalla por hito, elegida para que la diferencia con el hito anterior sea
legible sin repetir pesta\~nas de esclavo, temas o vistas auxiliares.

\newcommand{\spahito}[3]{%
\begin{figure}[p]
\centering
\includegraphics[width=0.88\textwidth,height=0.8\textheight,keepaspectratio]{esp_historico/#1}
\caption{Hito #2. #3}
\end{figure}
\clearpage}

\spahito{01_ui_inicial/01_vista_inicial_oscura.png}{1 --- interfaz inicial (08/06/2026)}{P\'agina \'unica sin pesta\~nas: ARM, START/STOP, gr\'afica con zoom, configuraci\'on de tres esclavos, exportaci\'on ZIP/CSV y log de WebSocket, todo visible a la vez. Es el punto de partida frente al que se miden los seis hitos siguientes.}

\spahito{02_funcional_incompleta/01_captura.png}{2 --- primera separaci\'on funcional (08/06/2026)}{Cambio principal: la p\'agina \'unica se reemplaza por pesta\~nas (Captura, Nodos, Esclavos, Export, Log). Define la arquitectura de navegaci\'on que se mantiene en todos los hitos posteriores; a\'un no agrega funciones de campo nuevas.}

\spahito{03_autocalibracion/01_captura.png}{3 --- autocalibraci\'on y lectura (15/06/2026)}{Cambio principal: aparecen cursores, vista espectral y un estado expl\'icito de autocalibraci\'on por esclavo. La interfaz deja de ser solo transporte y empieza a asistir la evaluaci\'on de la se\~nal.}

\spahito{04_pagina_operativa/01_captura.png}{4 --- p\'agina operativa (01/07/2026)}{Cambio principal: pesta\~nas Nodos y Export se consolidan en Esclavos y Preservados; se agregan autenticaci\'on y guardado CSV/ZIP integrado a la captura. La interfaz pasa de banco de pruebas a flujo de campo reconocible.}

\spahito{05_ui_finalizada/01_captura_completa.png}{5 --- captura finalizada (12/07/2026)}{Cambio principal: decimaci\'on por $N$, duraciones independientes para GEO y HAMMER, y l\'imites de captura por RAM y por SD. Estos controles provienen de restricciones reales de memoria detectadas durante las primeras pruebas de campo.}

\spahito{06_enlace_mdns/01_captura.png}{6 --- enlace e ingesta por mDNS (24/07/2026)}{Cambio principal: nueva pesta\~na Enlace con escaneo de redes de \SI{2.4}{\giga\hertz}, credenciales, URL de ingesta y cola de subida. Completa la ruta desde el maestro hacia una direcci\'on de servidor conocida, ausente en los cinco hitos anteriores.}

\spahito{07_pga_pgaout/01_captura_completa.png}{7 --- revisi\'on estable con PGA/PGAout (03/08/2026)}{Cambio principal: selecci\'on independiente entre PGA y PGAout por esclavo, sumada a rango de ADC y calibraci\'on ya estabilizados. Es la versi\'on vigente y la que aparece tambi\'en en la Fig.~\ref{fig:spa-hito7-resumen}.}


\else

\section*{Ap\'endice: recorrido visual exhaustivo de la SPA embebida}
\label{sec:anexo-spa-visual}
\addcontentsline{toc}{section}{Ap\'endice: recorrido visual exhaustivo de la SPA embebida}

Este ap\'endice re\'une, sin resumir, las capturas de los siete hitos de la SPA
mencionados en la Tabla~\ref{tab:spa-historica} y ya introducidos en la
Fig.~\ref{fig:spa-hito1-resumen} y la Fig.~\ref{fig:spa-hito7-resumen}. Se
mantiene como banco visual completo ---incluidas las pesta\~nas de configuraci\'on
que se repiten por nodo--- para que, al preparar una versi\'on resumida de la
presentaci\'on, puedan elegirse desde aqu\'i las pocas im\'agenes adicionales que
hagan falta, en lugar de volver a capturar pantallas. Existe tambi\'en una
versi\'on condensada de esta secci\'on, \texttt{cuarta-seccion-resumida.tex}, que
reemplaza este ap\'endice por una \'unica captura por hito.

El servidor de demostraci\'on emula WebSocket, estados, tres nodos, escaneo de
redes e ingesta local. Las formas de onda son \textbf{se\~nales sint\'eticas
utilizadas para activar y documentar la interfaz}; no se presentan como
mediciones de campo ni como validaci\'on del sistema.

\newcommand{\spavisual}[2]{%
\begin{figure}[p]
\centering
\includegraphics[width=\textwidth,height=0.79\textheight,keepaspectratio]{esp_historico/#1}
\caption{#2}
\end{figure}
\clearpage}

\clearpage
\paragraph{Hito 1 --- interfaz inicial del 08/06/2026.}
La primera versi\'on utiliza una p\'agina continua. Aunque todav\'ia no separa el
flujo en men\'us, ya re\'une adquisici\'on, gr\'aficas, configuraci\'on individual,
exportaci\'on y diagn\'ostico.

\spavisual{01_ui_inicial/01_vista_inicial_oscura.png}{Hito 1, vista oscura completa. La p\'agina continua muestra la primera integraci\'on de captura, gr\'aficas, nodos, configuraci\'on y registro.}
\spavisual{01_ui_inicial/02_vista_inicial_clara.png}{Hito 1, tema claro. La misma interfaz permite comparar contraste y legibilidad sin cambiar el flujo de operaci\'on.}
\spavisual{01_ui_inicial/03_captura_arm_start_stop.png}{Hito 1, bloque de adquisici\'on: cantidad de esclavos, ARM, duraci\'on, START, STOP y consulta de estado.}
\spavisual{01_ui_inicial/04_grafico_zoom_reset.png}{Hito 1, gr\'afica temporal y primeras herramientas de zoom y restablecimiento.}
\spavisual{01_ui_inicial/04b_nodos_activos.png}{Hito 1, detalle de los controles iniciales del filtro FIR, remoci\'on de continua, notch y pruebas de latencia.}
\spavisual{01_ui_inicial/05_descarga_zip.png}{Hito 1, exportaci\'on: nombre de la adquisici\'on, opci\'on CSV y descarga del contenedor ZIP.}
\spavisual{01_ui_inicial/06_log_websocket.png}{Hito 1, registro WebSocket. El log hace observable la conexi\'on y los mensajes intercambiados con el maestro.}
\spavisual{01_ui_inicial/07_esclavo_1_configuracion.png}{Hito 1, configuraci\'on del esclavo 1: tipo, VDAC, PGA, FIR, remoci\'on de continua, notch y comandos de prueba.}
\spavisual{01_ui_inicial/07_esclavo_2_configuracion.png}{Hito 1, configuraci\'on del esclavo 2. La repetici\'on por nodo anticipa el control distribuido posterior.}
\spavisual{01_ui_inicial/07_esclavo_3_configuracion.png}{Hito 1, configuraci\'on del esclavo 3 con los mismos controles y estad\'isticas independientes.}

\clearpage
\paragraph{Hito 2 --- primera separaci\'on funcional del 08/06/2026.}
La segunda revisi\'on reemplaza la p\'agina continua por pesta\~nas. La separaci\'on
todav\'ia es incompleta, pero define la arquitectura de navegaci\'on que se
mantendr\'a en los hitos siguientes.

\spavisual{02_funcional_incompleta/01_captura.png}{Hito 2, pesta\~na Captura con operaci\'on global y gr\'aficas visibles en paralelo.}
\spavisual{02_funcional_incompleta/02_nodos.png}{Hito 2, pesta\~na Nodos: estado compacto de los canales antes de entrar a su configuraci\'on detallada.}
\spavisual{02_funcional_incompleta/03_esclavos.png}{Hito 2, pesta\~na Esclavos con los paneles de configuraci\'on individual.}
\spavisual{02_funcional_incompleta/04_export.png}{Hito 2, pesta\~na Export: se independiza la generaci\'on de archivos del control de captura.}
\spavisual{02_funcional_incompleta/05_log.png}{Hito 2, pesta\~na Log con la secuencia de conexi\'on, estados y mensajes de los nodos.}
\spavisual{02_funcional_incompleta/06_tema_claro.png}{Hito 2, recorrido completo en tema claro.}

\clearpage
\paragraph{Hito 3 --- autocalibraci\'on y herramientas de lectura del 15/06/2026.}
En esta revisi\'on la interfaz deja de limitarse al transporte: aparecen cursores,
espectro, filtrado local y un estado expl\'icito de autocalibraci\'on.

\spavisual{03_autocalibracion/01_captura.png}{Hito 3, Captura: control global, frecuencia de muestreo y acceso directo a cursores y espectro.}
\spavisual{03_autocalibracion/02_nodos.png}{Hito 3, Nodos: lectura simult\'anea del estado y las se\~nales de los canales activos.}
\spavisual{03_autocalibracion/03_esclavos.png}{Hito 3, Esclavos: autocalibraci\'on, FIR, remoci\'on de continua y diagn\'ostico por nodo.}
\spavisual{03_autocalibracion/04_export.png}{Hito 3, Export: selecci\'on y preparaci\'on de los datos a conservar.}
\spavisual{03_autocalibracion/05_cursores.png}{Hito 3, cursores temporales incorporados a las gr\'aficas.}
\spavisual{03_autocalibracion/06_espectro.png}{Hito 3, vista espectral para inspeccionar la distribuci\'on de energ\'ia de la se\~nal.}
\spavisual{03_autocalibracion/07_autocalibracion.png}{Hito 3, autocalibraci\'on solicitada y filtro FIR activo en el panel del esclavo.}
\spavisual{03_autocalibracion/08_log.png}{Hito 3, registro de las operaciones de calibraci\'on, filtrado y comunicaci\'on.}
\spavisual{03_autocalibracion/09_tema_claro.png}{Hito 3, tema claro con todas las pesta\~nas y gr\'aficas.}

\clearpage
\paragraph{Hito 4 --- p\'agina operativa del 01/07/2026.}
El cuarto hito organiza un flujo de campo reconocible: adquisici\'on, configuraci\'on
de esclavos, preservaci\'on, exportaci\'on y acceso protegido.

\spavisual{04_pagina_operativa/01_captura.png}{Hito 4, Captura operativa con tres nodos y trazas sint\'eticas para ejercitar la visualizaci\'on.}
\spavisual{04_pagina_operativa/02_esclavos.png}{Hito 4, Esclavos: inversi\'on de se\~nal, PGA, FIR y funciones de hardware PSoC.}
\spavisual{04_pagina_operativa/03_preservados.png}{Hito 4, pesta\~na Preservados para separar adquisiciones conservadas del flujo vivo.}
\spavisual{04_pagina_operativa/04_preservar.png}{Hito 4, operaci\'on Preservar aplicada a la captura disponible.}
\spavisual{04_pagina_operativa/05_export_csv_zip.png}{Hito 4, nombre y guardado CSV/ZIP integrados en la pesta\~na de captura.}
\spavisual{04_pagina_operativa/06_autenticacion.png}{Hito 4, modal de acceso protegido emitido por el flujo real de autenticaci\'on WebSocket.}
\spavisual{04_pagina_operativa/07_log.png}{Hito 4, log operativo con estados, HELLO de nodos y confirmaciones del emulador.}
\spavisual{04_pagina_operativa/08_procesamiento_graficas.png}{Hito 4, envolventes y remoci\'on de componente continua aplicadas a las trazas de demostraci\'on.}
\spavisual{04_pagina_operativa/09_tema_claro.png}{Hito 4, tema claro de la p\'agina operativa.}

\clearpage
\paragraph{Hito 5 --- interfaz de captura finalizada del 12/07/2026.}
La revisi\'on agrega par\'ametros que provienen de restricciones reales de memoria
y de tipos de nodo: decimaci\'on, duraciones diferentes y l\'imites RAM/SD.

\spavisual{05_ui_finalizada/01_captura_completa.png}{Hito 5, vista completa del flujo de captura finalizado.}
\spavisual{05_ui_finalizada/02_decimacion_n.png}{Hito 5, aplicaci\'on del factor global de decimaci\'on $N$ y actualizaci\'on de la frecuencia efectiva.}
\spavisual{05_ui_finalizada/03_duraciones_geo_hammer.png}{Hito 5, duraciones independientes para nodos GEO y HAMMER.}
\spavisual{05_ui_finalizada/04_limite_ram.png}{Hito 5, acceso directo al m\'aximo de captura limitado por RAM.}
\spavisual{05_ui_finalizada/04b_limite_sd.png}{Hito 5, acceso directo al m\'aximo de captura habilitado por almacenamiento SD.}
\spavisual{05_ui_finalizada/05_descontar.png}{Hito 5, operaci\'on Descontar incorporada al control de la adquisici\'on.}
\spavisual{05_ui_finalizada/06_esclavos.png}{Hito 5, configuraci\'on avanzada: PGA, rango ADC, SD, FIR, rechazo de ruido de l\'inea y comandos PSoC.}
\spavisual{05_ui_finalizada/07_preservados.png}{Hito 5, selecci\'on de capturas preservadas antes de exportar o limpiar.}
\spavisual{05_ui_finalizada/08_log.png}{Hito 5, registro de comandos y estados de la interfaz finalizada.}
\spavisual{05_ui_finalizada/09_tema_claro.png}{Hito 5, tema claro.}
\spavisual{05_ui_finalizada/10_procesamiento_graficas.png}{Hito 5, controles de cruda, filtrada, envolventes y continua disponibles en la barra gr\'afica.}

\clearpage
\paragraph{Hito 6 --- enlace, ingesta y mDNS del 24/07/2026.}
Este hito completa la ruta desde el maestro hacia una direcci\'on conocida. La
SPA conserva la captura y suma escaneo, credenciales, diagn\'ostico de canal,
cola de archivos y subida del ZIP.

\spavisual{06_enlace_mdns/01_captura.png}{Hito 6, primera vista de la pesta\~na Enlace, con red, b\'usqueda, credenciales y URL del servidor reunidas en el maestro.}
\spavisual{06_enlace_mdns/02_esclavos.png}{Hito 6, Esclavos con RSSI ESP-NOW y configuraci\'on individual.}
\spavisual{06_enlace_mdns/03_preservados.png}{Hito 6, Preservados: capturas elegibles para el ZIP y la ingesta.}
\spavisual{06_enlace_mdns/04_enlace_estado.png}{Hito 6, Enlace: red, IP, canal, cantidad de archivos en cola y espacio libre.}
\spavisual{06_enlace_mdns/05_enlace_buscar.png}{Hito 6, b\'usqueda de redes de \SI{2.4}{\giga\hertz} visibles para el ESP32.}
\spavisual{06_enlace_mdns/06_enlace_configuracion.png}{Hito 6, configuraci\'on de SSID, credencial y URL de ingesta.}
\spavisual{06_enlace_mdns/07_subir_servidor.png}{Hito 6, operaci\'on Subir al server ejecutada contra la ruta local del emulador.}
\spavisual{06_enlace_mdns/08_log.png}{Hito 6, log de enlace, guardado de configuraci\'on y subida.}
\spavisual{06_enlace_mdns/09_tema_claro.png}{Hito 6, tema claro del flujo con conectividad.}

\clearpage
\paragraph{Hito 7 --- revisi\'on estable con PGA/PGAout del 03/08/2026.}
La \`ultima versi\'on mantiene el flujo anterior e incorpora el punto de selecci\'on
PGAout, junto con rango ADC, filtrado, calibraci\'on y lectura del estado de enlace.

\spavisual{07_pga_pgaout/01_captura_completa.png}{Hito 7, vista final completa con tres nodos y se\~nales sint\'eticas.}
\spavisual{07_pga_pgaout/02_grafico_cruda_filtrada.png}{Hito 7, controles preparados para alternar la visualizaci\'on temporal de la se\~nal cruda y filtrada.}
\spavisual{07_pga_pgaout/03_envolventes_dc.png}{Hito 7, se\~nal, envolvente y remoci\'on de continua despu\'es de actualizar el filtro.}
\spavisual{07_pga_pgaout/04_espectro.png}{Hito 7, espectro de los canales activos.}
\spavisual{07_pga_pgaout/05_cursores_zoom.png}{Hito 7, cursores y zoom sobre la vista temporal.}
\spavisual{07_pga_pgaout/06_captura_parametros.png}{Hito 7, ARM, decimaci\'on, duraciones GEO/HAMMER y l\'imites de memoria.}
\spavisual{07_pga_pgaout/07_esclavos_general.png}{Hito 7, vista general de los tres paneles de esclavo.}
\spavisual{07_pga_pgaout/07b_esclavo_filtro_adc_calibracion.png}{Hito 7, combinaci\'on de rango ADC, FIR pasa banda y calibraci\'on en un nodo GEO.}
\spavisual{07_pga_pgaout/08_esclavo_pga.png}{Hito 7, selecci\'on de ganancia PGA.}
\spavisual{07_pga_pgaout/09_esclavo_pgaout.png}{Hito 7, selecci\'on independiente de PGAout en el nuevo punto de la cadena anal\'ogica.}
\spavisual{07_pga_pgaout/10_preservados.png}{Hito 7, capturas preservadas en la versi\'on final.}
\spavisual{07_pga_pgaout/11_enlace_estado.png}{Hito 7, estado de red, canal y cola en la pesta\~na Enlace.}
\spavisual{07_pga_pgaout/12_enlace_escaneo.png}{Hito 7, resultado del escaneo de redes de \SI{2.4}{\giga\hertz}.}
\spavisual{07_pga_pgaout/13_enlace_servidor.png}{Hito 7, URL local de ingesta guardada en el maestro emulado.}
\spavisual{07_pga_pgaout/14_subida_servidor.png}{Hito 7, ZIP de demostraci\'on subido correctamente al endpoint local.}
\spavisual{07_pga_pgaout/15_log.png}{Hito 7, registro final de conexi\'on, HELLO, comandos, configuraci\'on y subida.}
\spavisual{07_pga_pgaout/16_tema_claro.png}{Hito 7, vista final en tema claro.}

\clearpage


\fi
